\begin{ouverture}[calcul]{2}{Repérer, comparer et arrondir}{Partie 1 · Nombres et calculs}
{Demi-droite · Ordre · Encadrement · Arrondi}
\objectif{Lire et placer une abscisse décimale sur une demi-droite graduée.}
\objectif{Comparer, ranger, encadrer et intercaler des nombres décimaux.}
\objectif{Donner l’arrondi d’un nombre à l’unité, au dixième ou au centième.}
\prerequis{Valeur des chiffres selon leur rang ; lecture d’une graduation.}
\end{ouverture}
\pagestyle{otMain}\essentiel\fontsize{11}{15}\selectfont
\phantomsection\label{cours-2-1}
\essSection{1}{Lire et construire une graduation}
\begin{definition}Sur une demi-droite graduée, l’origine est repérée par 0 et une même longueur, l’unité, est reportée régulièrement. Chaque point est repéré par un nombre, appelé son abscisse.\end{definition}
\label{fig-cours-2-1-cours-droite-zoom}\figureGeo{cours-droite-zoom}

\begin{savoirfaire}[Lire une graduation]Entre 3 et 4, il y a dix intervalles égaux : chacun vaut $0{,}1$. Entre $3{,}4$ et $3{,}5$, dix intervalles encore : chacun vaut $0{,}01$. On cherche toujours la valeur d’un petit intervalle avant de lire.\end{savoirfaire}
\essSection{2}{Comparer et ordonner des décimaux}
\begin{propriete}Pour comparer deux nombres décimaux, on compare les parties entières, puis les dixièmes, puis les centièmes… $12{,}7>12{,}69$, car 7 dixièmes est plus que 6 dixièmes. Écrire le même nombre de décimales aide : $12{,}70>12{,}69$.\end{propriete}
\begin{attention}Le nombre de chiffres ne décide pas : $0{,}8>0{,}75$ alors que 75 est plus grand que 8. On compare rang par rang, en commençant par la gauche.\end{attention}
\clearpage
\phantomsection\label{cours-2-2}
\essSection{3}{Encadrer et intercaler}
\begin{definition}Encadrer un nombre, c’est le placer entre deux nombres. $3<3{,}47<4$ est un encadrement à l’unité ; $3{,}4<3{,}47<3{,}5$ est un encadrement au dixième.\end{definition}
\begin{savoirfaire}[Intercaler]Entre $2{,}3$ et $2{,}4$, écrire $2{,}30$ et $2{,}40$ fait apparaître d’autres nombres : $2{,}31$, $2{,}35$… Entre deux décimaux différents, on peut toujours en intercaler un autre.\end{savoirfaire}
\essSection{4}{Arrondir et approcher un nombre}
\begin{definition}L’arrondi à l’unité d’un nombre est l’entier le plus proche ; l’arrondi au dixième, le nombre à une décimale le plus proche ; de même au centième. Si le chiffre qui suit le rang voulu est 5, 6, 7, 8 ou 9, on arrondit au-dessus ; sinon, au-dessous.\end{definition}
\label{fig-cours-2-2-cours-arrondi}\figureGeo{cours-arrondi}

\begin{exemple}$7{,}364$ s’arrondit à 7 à l’unité, à $7{,}4$ au dixième et à $7{,}36$ au centième. Certains nombres ne s’écrivent pas avec un nombre fini de décimales, comme $\frac13=0{,}333\ldots$ : on utilise alors une valeur arrondie, $0{,}33$ au centième.\end{exemple}

\clearpage\rubrique{Découverte}
\begin{blocLivre}{ex-2-1}\begin{exercice}[1]{Lire des abscisses}
La demi-droite est graduée en dixièmes. Lire les abscisses des points A, B et C.
\label{fig-ex-2-1}\figureGeo{rep-1}
\end{exercice}\end{blocLivre}
\begin{blocLivre}{ex-2-2}\begin{exercice}[1]{Placer des points}
La demi-droite est graduée en centièmes, de $5{,}2$ à $5{,}4$. Placer les points D, E et F d’abscisses $5{,}25$ ; $5{,}38$ et $5{,}3$.
\label{fig-ex-2-2}\figureGeo{rep-2}
\end{exercice}\end{blocLivre}
\begin{blocLivre}{ex-2-3}\begin{exercice}[1]{Comparer}
Comparer avec $<$ ou $>$ : $12{,}7$ et $12{,}69$ ; $0{,}8$ et $0{,}75$ ; $3{,}05$ et $3{,}5$ ; $7$ et $6{,}999$.
\end{exercice}\end{blocLivre}
\begin{blocLivre}{ex-2-4}\begin{exercice}[1]{Ranger}
Ranger dans l’ordre croissant : $4{,}2$ ; $4{,}02$ ; $4{,}25$ ; $4{,}205$ ; $4{,}3$.
\end{exercice}\end{blocLivre}
\begin{blocLivre}{ex-2-5}\begin{exercice}[1]{Arrondir}
Donner l’arrondi à l’unité, au dixième et au centième de $7{,}364$, puis de $12{,}58$.
\end{exercice}\end{blocLivre}

\rubrique{Entraînement}
\begin{blocLivre}{ex-2-6}\begin{exercice}[2]{Encadrer}
Encadrer $15{,}738$ à l’unité, au dixième, puis au centième.
\end{exercice}\end{blocLivre}
\begin{blocLivre}{ex-2-7}\begin{exercice}[2]{Intercaler}
Intercaler un nombre entre $2{,}3$ et $2{,}4$ ; entre $0{,}59$ et $0{,}6$ ; entre 8 et $8{,}01$.
\end{exercice}\end{blocLivre}
\begin{blocLivre}{ex-2-8}\begin{exercice}[2]{Zoom sur la droite}
La demi-droite est graduée en centièmes entre $1{,}4$ et $1{,}5$. Lire l’abscisse de G. Le point H est au milieu de deux graduations : donner son abscisse.
\label{fig-ex-2-8}\figureGeo{rep-8}
\end{exercice}\end{blocLivre}
\begin{blocLivre}{ex-2-9}\begin{exercice}[2]{Valeurs approchées}
On sait que $\frac13=0{,}333\ldots$ et $\frac23=0{,}666\ldots$, les chiffres se répétant sans fin. Donner les arrondis au dixième et au centième de ces deux nombres.
\end{exercice}\end{blocLivre}
\begin{blocLivre}{ex-2-10}\begin{exercice}[2]{Tous les candidats}
Trouver tous les nombres à deux décimales dont l’arrondi au dixième est $3{,}5$. Combien y en a-t-il ?
\end{exercice}\end{blocLivre}

\rubrique{Problèmes}
\begin{blocLivre}{ex-2-11}\begin{exercice}[3]{Le saut en longueur}
Résultats d’un concours de saut : Léna $4{,}85$ m ; Malik $4{,}9$ m ; Sofia $4{,}08$ m ; Tom $4{,}89$ m ; Inès $4{,}8$ m. Établir le classement et le podium.
\end{exercice}\end{blocLivre}
\begin{blocLivre}{ex-2-12}\begin{exercice}[3]{Les prix arrondis}
Au marché, Yanis achète trois articles à $2{,}49$ €, $13{,}95$ € et $0{,}99$ €. Arrondir chaque prix à l’euro, estimer le total, puis calculer le prix exact.
\end{exercice}\end{blocLivre}
\begin{blocLivre}{ex-2-13}\begin{exercice}[3]{Le chronomètre}
Quatre coureurs ont été chronométrés : Ana $12{,}48$ s ; Bilal $12{,}4$ s ; Chloé $12{,}52$ s ; Dan $12{,}405$ s. Qui gagne ? Si le chronomètre n’affichait que les dixièmes (arrondis), pourrait-on départager tout le monde ?
\end{exercice}\end{blocLivre}
\begin{blocLivre}{ex-2-14}\begin{exercice}[3]{Le nombre mystère}
Je suis un nombre à deux décimales. Mon arrondi à l’unité est 8, mon arrondi au dixième est $7{,}6$ et mon chiffre des centièmes est 9. Qui suis-je ?
\end{exercice}\end{blocLivre}
\begin{blocLivre}{ex-2-15}\begin{exercice}[3]{Vrai ou faux}
Hugo affirme : « $3{,}15$ est plus grand que $3{,}9$, car 15 est plus grand que 9. » Corriger son raisonnement.
\end{exercice}\end{blocLivre}

\clearpage
